\documentclass[11pt]{article}
\setlength{\topmargin}{-.3in}
\setlength{\oddsidemargin}{-0.1in}
\setlength{\evensidemargin}{-0.1in}
\setlength{\textwidth}{6.7in}
\setlength{\headheight}{0in}
\setlength{\headsep}{0in}
\setlength{\topskip}{0.5in}
\setlength{\textheight}{9.25in}
\setlength{\parindent}{0.0in}
\setlength{\parskip}{1em}
\usepackage{macros_cavag}
\usepackage{graphicx}
\usepackage{hyperref}
\usepackage{epsfig}

\pagestyle{empty}

\begin{document}

\today

Dr. G. Fabbiano\\
Smithsonian Astrophysical Observatory\\
60 Garden Street\\
MS 6\\
Cambridge, MA 02138\\

Dear Dr. Fabbiano:

Please accept the attached application for your postdoctoral
position advertised in the November 2007 issue of the AAS Job
Register. For my thesis (advised by Megan Donahue and  Mark Voit),
entitled {\it `Virialization, Entropy, and Feedback in Clusters of
Galaxies'}, I am studying the coupling of feedback mechanisms -- such
as AGN, star formation, and conduction in cluster cores -- to gas
entropy, and the role of this feedback in altering global ICM
properties and truncating the high mass end of the galaxy luminosity
function. I have also been studying a method for quantifying the
virialization state of clusters through the band dependence of X-ray
temperatures. For my thesis I assembled a sample of 350 archival
{\textit{Chandra}} observations for 276 clusters totaling 11.6 Msec of
data. The results of this laborious effort have been many and are
detailed in my research summary.

I am a great research asset for any X-ray astronomer because of my
extensive experience with X-ray data, innovative technical skills, and
ability to independently advance the group's and my own research
objectives. I feel the post-doctoral position under your advisory at
SAO is an excellent fit for me, and your research goals will benefit
from my addition. My expertise in X-ray astronomy ideally suits me to
further work on better understanding AGN, models of galaxy formation,
and the effect of feedback in shaping galaxy and cluster
properties. As our understanding of feedback mechanisms in clusters
improves, additional focus must be paid to these same processes in
groups and individual galaxies if we are to have a complete
understanding of hierarchical structure formation.

Along with this letter are my resum\'{e}, a summary of past and
current research, and a brief description of possible research
directions. Letters of recommendation from Megan Donahue, Mark Voit,
and Jack Baldwin are available at anytime via their contact
information provided in my resum\'{e}. Please contact me if there is
any further information I can provide as you review my application.

Thank you for your consideration.

Sincerely,\\
\begin{minipage}{7.5in}
\epsfig{file=signature.ps,height=0.75in}
\end{minipage}
Kenneth W. Cavagnolo\\
Michigan State University

\end{document}
